% SPDX-License-Identifier: LPPL-1.3c
% Copyright (C) 2026 Christophe Jorssen
% Author and Current Maintainer: Christophe Jorssen <christophe.jorssen@gmail.com>
% LPPL maintenance status: maintained
% This file is part of luafemm. See LICENSE and LICENSES.md for its terms.
\input tikz.tex
\usetikzlibrary{femm}
\nopagenumbers
\tikzset{test box/.style={femm/problem={xmin=-10,xmax=10,ymin=-10,ymax=10,mesh size=4}},
  imposed field/.style={femm/boundary={y coefficient=#1}},
  imposed field/.default=.2,
  imposed field/.append style={femm/boundary={side=left,type=neumann}}}
% Parameterised styles, append style, and multiple problem keys. The begin
% hook after the first problem key must already see the complete setup.
\tikzpicture[imposed field,test box,
  execute at begin picture={\directlua{
    assert(luafemm.current.boundaries.left.type=='neumann')
    assert(luafemm.current.boundaries.top.dy==.2)
  }},femm/problem={mesh size=3}]
  \femmsolve
  \directlua{luafemm.boundary_reference=luafemm.current}
\endtikzpicture
% Key setup and command setup must give exactly the same discrete result.
\tikzpicture[test box,femm/problem={mesh size=3}]
  \femmboundary[y coefficient=.2]
  \femmboundary[side=left,type=neumann]
  \femmsolve
  \directlua{
    local old=luafemm.boundary_reference
    local m=luafemm.current
    assert(rawlen(old.A)==rawlen(m.A))
    for i,a in ipairs(old.A) do assert(a==m.A[i]) end
  }
\endtikzpicture
% Defaults declared outside a picture apply independently to every new model.
\begingroup
\tikzset{femm/boundary={y coefficient=.1}}
\tikzpicture[test box]
  \directlua{assert(luafemm.current.boundaries.top.dy==.1)}
  \scope[femm/boundary={side=left,type=neumann}]
  \endscope
  % A physical declaration survives the graphical scope that applied it.
  \directlua{assert(luafemm.current.boundaries.left.type=='neumann')}
  \path[femm/boundary={side=right,type=neumann}] (0,0)--(1,0);
  \directlua{assert(luafemm.current.boundaries.right.type=='neumann')}
  \femmsolve
\endtikzpicture
\tikzpicture[test box]
  \directlua{
    assert(luafemm.current.boundaries.top.dy==.1)
    assert(luafemm.current.boundaries.left.type=='dirichlet')
    assert(luafemm.current.boundaries.right.type=='dirichlet')
  }
\endtikzpicture
\endgroup
% An every-picture style is inherited normally, including its parameter.
\begingroup
\tikzset{every picture/.append style={imposed field=.4}}
\tikzpicture[test box]
  \directlua{
    assert(luafemm.current.boundaries.top.dy==.4)
    assert(luafemm.current.boundaries.left.type=='neumann')
  }
\endtikzpicture
\endgroup
% Neither pending declarations nor the ready flag leak between pictures.
\tikzpicture[test box]
  \directlua{
    for _,p in pairs(luafemm.current.boundaries) do
      assert(p.type=='dirichlet' and p.dy==0)
    end
  }
  \tikzset{femm/boundary={y coefficient=.7},femm/boundary}
  \femmsolve
  \directlua{assert(luafemm.current.stats.peak==0)}
\endtikzpicture
% Exercise the actual TeX key after meshing. The expected Lua rejection is
% intercepted at the bridge so that this positive compilation can assert it.
\tikzpicture[test box]
  \femmsolve
  \directlua{
    luafemm.saved_boundary_bridge=luafemm.tex.boundary
    luafemm.tex.boundary=function(...)
      local ok,message=pcall(luafemm.saved_boundary_bridge,...)
      assert(not ok and tostring(message):find('declare boundary conditions before meshing',1,true))
      luafemm.boundary_rejection_seen=true
    end
  }
  \tikzset{femm/boundary={side=left,type=neumann}}
  \directlua{
    luafemm.tex.boundary=luafemm.saved_boundary_bridge
    luafemm.saved_boundary_bridge=nil
    assert(luafemm.boundary_rejection_seen)
    assert(luafemm.current.boundaries.left.type=='dirichlet')
  }
\endtikzpicture
\directlua{texio.write_nl('PASS: boundary keys, styles, timing, grouping and command equivalence')}
\bye
